% sections/06-conclusion.tex

\section{Conclusion}\label{sec:conclusion}

We have introduced eight energy functionals for the Kempe chain
reconfiguration graph (Section~\ref{sec:methods}), grounded in the
tetrahedral color embedding. These functionals provide complementary
views of the reconfiguration process, connecting graph coloring to
anti-ferromagnetic spin models, polyhedral geometry, and discrete
potential theory.

\paragraph{Principal findings.}
Evaluation on the counterexamples at $n = 9$
(Section~\ref{sec:results}) identified two discriminating signatures
that persist at $n = 10$: (1)~local entropy at the critical vertex is
systematically higher for unsafe swaps, and (2)~the Magic Gem energy
landscape exhibits a two-basin kinetic trap structure separating safe
from unsafe paths. A third candidate signal — surface tension rigidity
correlating with merge-proneness — held at $n = 9$ but was
subsequently falsified at $n \geq 10$
(Remark~\ref{rem:st-falsified}).

\paragraph{The Surface Tension Rigidity Conjecture.}
We formalized finding~(3) as
Conjecture~\ref{conj:st-rigidity}: zero surface tension rigidity is
a sufficient condition for a Kempe swap to be unsafe. This conjecture
held perfectly at $n = 9$, and we speculated about connections to sheaf
cohomology and the Penrose evaluation
(Section~\ref{sec:discussion}), though these remain unverified.
Subsequent computation at $n \geq 10$ falsified the conjecture: the
correlation between zero rigidity and merge-proneness breaks down at
larger scales (see Remark~\ref{rem:st-falsified}). The local entropy
discrimination and kinetic trap structure, by contrast, persist in the
extended dataset.

\paragraph{Open questions.}
\begin{itemize}
  \item \textbf{Persistence:} Do the local entropy discrimination and
    kinetic trap structure persist at all scales? Enumeration at
    $n = 10$ (${\sim}1.9$ million cases) is consistent with
    persistence, but the Surface Tension Rigidity Conjecture fails at
    this scale (Remark~\ref{rem:st-falsified}).
  \item \textbf{Refined indicators:} Can surface tension be combined
    with local entropy or Magic Gem energy into a composite indicator
    that persists where surface tension alone does not?
  \item \textbf{TQFT:} Can the suggestive analogy between
    Magic Gem energy basins and Penrose evaluation signs
    (Section~\ref{sec:tqft}) be made precise and extended beyond
    $n = 9$? Even a positive answer would require substantial
    additional work to yield a rigorous proof of Safe Path Existence.
  \item \textbf{Composite functionals:} Can the eight functionals be
    combined into a single ``reconfiguration potential'' whose gradient
    flow always finds safe paths? Constructing such a potential and
    proving it works would be a major undertaking beyond the scope of
    the present descriptive analysis.
\end{itemize}

\paragraph{Broader perspective.}
The energy functionals introduced here are descriptive tools
(Section~\ref{sec:limitations}): they reveal structure in the
reconfiguration landscape but do not, by themselves, prove that safe
paths exist. Two of the three discriminating signatures (local entropy
and kinetic traps) persist at $n = 10$, while the third (surface
tension rigidity) does not. The connections to TQFT, sheaf cohomology,
chromatic polynomials, and discharging remain speculative. Whether these
analogies can be developed into rigorous proof techniques is an open
question that the present work does not resolve.

\paragraph{Reproducibility.}
All computational results reported in this paper can be reproduced from
the accompanying code repository. Triangulation counts at each order
were independently verified against OEIS sequence
A000109~\cite{oeis_A000109}. The enumeration, energy computation, and
safe-path search pipelines are documented in the repository README.
